% ============ §2 名词 ============
\hdr{§2}{名词：先定性，再定数}

\callout{\textbf{这一节的坐标}：У4 p9--p19（кто／что、三个性、-ь 的两栏词表、Читайте 练习）、У5 p2--p14（复数的三个表、只有单数／只有复数的词）、У7 p17--p19（特殊复数与只有复数的词，同一批规则在 У7 又出现一次）。\\
\textbf{这一节要解决的事}：给一个名词，你能判出它的性、写出它的复数，并说出为什么是这个词尾。}

\subhdr{2.1　名词只回答两个问题：Кто? 和 Что?}

У4 p9 把名词的入口收成两个疑问词：\textbf{Кто?}（谁）问人和动物，\textbf{Что?}（什么）问东西和事情。答句里更常用的是代词而不是重复名词（У4 p13、p17）：

\begin{minipage}{\linewidth}
\tabtop
\begin{tabular}{@{}>{\raggedright\arraybackslash}p{0.14\linewidth}>{\raggedright\arraybackslash}p{0.26\linewidth}>{\raggedright\arraybackslash}p{0.50\linewidth}@{}}
\thead{性} & \thead{代词} & \thead{课件例词} \\
\end{tabular}
\tabmid
\begin{tabular}{@{}>{\raggedright\arraybackslash}p{0.14\linewidth}>{\raggedright\arraybackslash}p{0.26\linewidth}>{\raggedright\arraybackslash}p{0.50\linewidth}@{}}
阳性 & он & Ива́н, дом, парк, заво́д, музе́й, слова́рь \\
阴性 & она́ & А́нна, ко́мната, ла́мпа, пе́сня, аудито́рия, тетра́дь \\
中性 & оно́ & окно́, молоко́, мо́ре, зда́ние, вре́мя \\
\end{tabular}
\tabbot
\end{minipage}

\textbf{这张表在说什么}：它把「性」从抽象语法变成一件具体的事——\textbf{你用一个什么代词去接它}。因此「这个词是什么性」等价于「这个词能不能用 он 顶替」。У4 p17 的句子就是这个用法：\textbf{Где соба́ка? — Она́ до́ма.}（соба́ка 以 -а 结尾 → она́）、\textbf{Когда́ уро́к? — Он за́втра.}（уро́к 以辅音结尾 → он）。

\textbf{逐行读法}：阳性那一行几乎都以辅音结尾（дом、парк、заво́д）， музе́й 以 -й 结尾、слова́рь 以 -ь 结尾，这两个是后面 2.2、2.3 的重点；阴性那一行以 -а／-я 为主，最后一个 тетра́дь 是 -ь；中性以 -о／-е 为主，вре́мя 是 -мя。

\subhdr{2.2　三个性与词尾：能推的推，推不出的记}

У4 p11--p12 的原表只给了词尾，没有给「为什么」。先把原表抄清：

\begin{minipage}{\linewidth}
\tabtop
\begin{tabular}{@{}>{\raggedright\arraybackslash}p{0.12\linewidth}>{\raggedright\arraybackslash}p{0.10\linewidth}>{\raggedright\arraybackslash}p{0.26\linewidth}>{\raggedright\arraybackslash}p{0.38\linewidth}@{}}
\thead{性} & \thead{代词} & \thead{词尾} & \thead{例词} \\
\end{tabular}
\tabmid
\begin{tabular}{@{}>{\raggedright\arraybackslash}p{0.12\linewidth}>{\raggedright\arraybackslash}p{0.10\linewidth}>{\raggedright\arraybackslash}p{0.26\linewidth}>{\raggedright\arraybackslash}p{0.38\linewidth}@{}}
阳性 & он & 硬辅音（□）、-й、-ь ※ & дом, музе́й, слова́рь \\
阴性 & она́ & -а、-я、-ь ※ & ко́мната, пе́сня, тетра́дь \\
中性 & оно́ & -о、-е、-мя & окно́, мо́ре, вре́мя \\
特例 & он & -а／-я 却是阳性 & па́па, де́душка, дя́дя, мужчи́на, ю́ноша \\
\end{tabular}
\tabbot
\end{minipage}

\textbf{这张表在说什么}：三行里有两行是干净的（阴性看 -а／-я，中性看 -о／-е），有一行是脏的——\textbf{-ь 同时出现在阳性和阴性两行}，课件在那里打了感叹号，意思是「光看词尾判不出来」。

\textbf{逐行读法}：第一行「硬辅音」指词尾是 б в г д ж з к л м н п р с т ф х ц ч ш щ 之一（дом、стол、парк）；第二行 -а／-я 是阴性的主力（ко́мната、пе́сня）；第三行 -о／-е 是中性主力（окно́、мо́ре），另有极少几个 -мя（вре́мя、и́мя）。

\textbf{补：-ь 结尾怎么办。}课件给的是两栏词表（下一节），没有给判断法。本讲义按俄语构词常识补四条可推的规律，\textbf{这四条是归纳，不是课件原文}：

\begin{itemize}
\item \textbf{ж／ш／ч／щ ＋ ь → 阴性}：ночь、по́лночь、вещь、дочь。因为阳性名词在这四个字母后面\textbf{不带 ь}（нож、мяч、врач、плащ）——所以凡是在它们后面见到 ь，一定是阴性。
\item \textbf{-ость 结尾 → 阴性}：но́вость。
\item \textbf{-тель 结尾 → 阳性}：преподава́тель。
\item \textbf{-арь 结尾 → 阳性}：слова́рь、календа́рь。
\end{itemize}

四条都不符的，只能查表（\S2.3 那 38 个词）。

\warnbox{\textbf{边界一}：特例只限「指人的 -а／-я」，而且数量很少：\textbf{па́па、де́душка、дя́дя、мужчи́на、ю́ноша}。别把这条推广成「-а 结尾也可能是阳性」——那样会写错 99\% 的词。\\
\textbf{边界二}：外来不变格词（\textbf{ко́фе、метро́、пальто́、такси́}）的性要单独记：ко́фе 是阳性（он），其余多是中性（оно́），详见 \S5.4。}

\subhdr{2.3　-ь 的 38 个词：课件 p13 的两栏照抄}

У4 p13 把 -ь 结尾的名词分成左右两栏，各 19 个。这张表没有规律可压缩，只能整词记——下面照抄，并给出中文：

\begin{minipage}{\linewidth}
\tabtop
\begin{tabular}{@{}>{\raggedright\arraybackslash}p{0.23\linewidth}>{\raggedright\arraybackslash}p{0.23\linewidth}>{\raggedright\arraybackslash}p{0.23\linewidth}>{\raggedright\arraybackslash}p{0.22\linewidth}@{}}
\thead{阳性：он} & \thead{中文} & \thead{阴性：она́} & \thead{中文} \\
\end{tabular}
\tabmid
\begin{tabular}{@{}>{\raggedright\arraybackslash}p{0.23\linewidth}>{\raggedright\arraybackslash}p{0.23\linewidth}>{\raggedright\arraybackslash}p{0.23\linewidth}>{\raggedright\arraybackslash}p{0.22\linewidth}@{}}
слова́рь & 词典 & тетра́дь & 练习本 \\
преподава́тель & 教师 & пло́щадь & 广场；面积 \\
день & 白天、日子 & ночь & 夜 \\
календа́рь & 日历 & дверь & 门 \\
портфе́ль & 公文包 & о́сень & 秋天 \\
карто́фель & 土豆 & крова́ть & 床 \\
автомоби́ль & 汽车 & но́вость & 新闻 \\
дождь & 雨 & ме́бель & 家具 \\
ка́шель & 咳嗽 & о́бувь & 鞋类 \\
кора́бль & 船 & любо́вь & 爱 \\
ли́вень & 暴雨 & ладо́нь & 手掌 \\
ло́коть & 肘 & мете́ль & 暴风雪 \\
но́готь & 指甲 & мысль & 思想 \\
ноль & 零 & о́чередь & 队列；轮次 \\
ого́нь & 火 & па́мять & 记忆 \\
по́лдень & 正午 & посте́ль & 床铺 \\
реме́нь & 皮带 & цель & 目标 \\
юа́нь & 元（钱） & по́лночь & 午夜 \\
спекта́кль & 戏剧 & меда́ль & 奖章 \\
\end{tabular}
\tabbot
\end{minipage}

\textbf{这张表在说什么}：38 个词里，阳性 19、阴性 19，正好一半一半——这就是「-ь 判不出来」的实证。左栏多是与「人、器物、时间、职业」相关的词，右栏多与「抽象概念、女性相关物品、时间」相关，但这个倾向很弱，\textbf{不要拿来当规则用}。

\textbf{逐行读法}：最容易混的四对已经在同一行里：день（阳）／ночь（阴）、дождь（阳）／тетра́дь（阴）、портфе́ль（阳）／посте́ль（阴）、юа́нь（阳）／меда́ль（阴）。背的时候成对背，比单背记得牢。

\aside{本节四条归纳规律（ж／ш／ч／щ＋ь、-ость、-тель、-арь）能覆盖的是「有标记」的词；这 38 个里有相当一部分没有标记，所以必须整词记。判断顺序：\textbf{先套四条规律 → 四条不符就查这张表}。}

\subhdr{2.4　复数：一条规律，三个性别各读一遍}

У5 p6--p8 的三张表（阳性、阴性、中性）表面上是三套词尾，其实是同一条规律在三个性别上的读数：

\begin{minipage}{\linewidth}
\tabtop
\begin{tabular}{@{}>{\raggedright\arraybackslash}p{0.12\linewidth}>{\raggedright\arraybackslash}p{0.22\linewidth}>{\raggedright\arraybackslash}p{0.18\linewidth}>{\raggedright\arraybackslash}p{0.40\linewidth}@{}}
\thead{性} & \thead{单数词尾} & \thead{复数词尾} & \thead{课件例词} \\
\end{tabular}
\tabmid
\begin{tabular}{@{}>{\raggedright\arraybackslash}p{0.12\linewidth}>{\raggedright\arraybackslash}p{0.22\linewidth}>{\raggedright\arraybackslash}p{0.18\linewidth}>{\raggedright\arraybackslash}p{0.40\linewidth}@{}}
阳 & 硬辅音 & -ы & студе́нт → студе́нты；университе́т → университе́ты \\
阳 & г／к／х／ж／ш／ч／щ & -и & уро́к → уро́ки；нож → но́жи；каранда́ш → карандаши́ \\
阳 & -ь、-й & -и & слова́рь → словари́；музе́й → музе́и \\
阴 & -а & -ы & сестра́ → сёстры；газе́та → газе́ты \\
阴 & -га／-ка／-ха／-жа／-ша／-ча／-ща & -и & кни́га → кни́ги；му́ха → му́хи \\
阴 & -ь、-я & -и & тетра́дь → тетра́ди；пе́сня → пе́сни \\
中 & -о & -а & окно́ → о́кна \\
中 & -е、-ие & -я、-ия & по́ле → поля́；упражне́ние → упражне́ния \\
\end{tabular}
\tabbot
\end{minipage}

\textbf{这张表在说什么}：八个字——\textbf{硬词干加 ы／а，软词干加 и／я}。「软词干」指词干末尾是软辅音（-ь、-й、-я、-е 结尾的词都算）。所以阳性硬辅音用 -ы、г／к／х／ж／ш／ч／щ 后用 -и（这七个字母后面永远不写 ы，见 \S4.2）、-ь／-й 后用 -и；阴性 -а 用 -ы、-га 等用 -и；中性硬词干 -о 用 -а、软词干 -е 用 -я。

\textbf{逐行读法}：第 2 行和第 5 行是同一个规则的两个性别版本（七个字母后写 и）；第 3 行和第 6 行也是同一规则（-ь 一律 -и）；第 1、7 行是「硬 → ы／а」；第 8 行是「软 → я／ия」。

\textbf{表后要提醒的一个坑}：\textbf{студе́нт → студе́нты} 但 \textbf{студе́нтка → студе́нтки}，因为 -тка 的 к 在七个字母里；同理 \textbf{сестра́ → сёстры} 里 е 变成 ё，是因为重音跑到了 ё 上（сёстры）——这类「重音跑位」只能遇到再记，没有规则。

\subhdr{2.5　规则之外的：24 组特殊复数与「只有一种数」的词}

У7 p18 单独列了一批「复数不按规则来」的词，У5 p9 与 У7 p19 又列了两批「只有单数／只有复数」的词：

\begin{minipage}{\linewidth}
\tabtop
\begin{tabular}{@{}>{\raggedright\arraybackslash}p{0.20\linewidth}>{\raggedright\arraybackslash}p{0.24\linewidth}>{\raggedright\arraybackslash}p{0.20\linewidth}>{\raggedright\arraybackslash}p{0.24\linewidth}@{}}
\thead{单数} & \thead{复数} & \thead{单数} & \thead{复数} \\
\end{tabular}
\tabmid
\begin{tabular}{@{}>{\raggedright\arraybackslash}p{0.20\linewidth}>{\raggedright\arraybackslash}p{0.24\linewidth}>{\raggedright\arraybackslash}p{0.20\linewidth}>{\raggedright\arraybackslash}p{0.24\linewidth}@{}}
мать & ма́тери & дочь & до́чери \\
сестра́ & сёстры & оте́ц & отцы́ \\
день & дни & и́мя & имена́ \\
вре́мя & времена́ & я́блоко & я́блоки \\
ребёнок & де́ти & челове́к & лю́ди \\
дом & дома́ & го́род & города́ \\
бе́рег & берега́ & лес & леса́ \\
о́стров & острова́ & по́езд & поезда́ \\
ве́чер & вечера́ & брат & бра́тья \\
сын & сыновья́ & друг & друзья́ \\
муж & мужья́ & де́рево & дере́вья \\
лист & ли́стья & стул & сту́лья \\
\end{tabular}
\tabbot
\end{minipage}

\textbf{这张表在说什么}：24 组里其实藏着三类不同的「不规则」——第一类\textbf{词干变了}（мать → ма́тери、и́мя → имена́、ребёнок → де́ти、челове́к → лю́ди）；第二类\textbf{重音跑到词尾}（дом → дома́、го́род → города́、лес → леса́、ве́чер → вечера́）；第三类\textbf{加了 -ья- 一类的后缀}（брат → бра́тья、стул → сту́лья、де́рево → дере́вья）。

\textbf{逐行读法}：读这张表时要连着\textbf{重音}一起读——第二类（дом、го́род、лес、ве́чер）单数是词干重音，复数重音跑到最后一个音节；第三类（брат、стул、де́рево、лист）复数多出一个 -ья-。

\textbf{只有单数}（У5 p9）：со́лнце、луна́、молоко́、мя́со、рис、карто́фель、лапша́、лук。
\textbf{只有复数}（У7 p19）：роди́тели、де́ньги、часы́、очки́、но́жницы、брю́ки、ша́хматы、джи́нсы、духи́。

\warnbox{\textbf{边界}：「只有复数」不等于「复数形式难写」，而是\textbf{它没有单数可用}。часы́ 指「钟／表」这一整件东西，说「一只表」不用单数，而说 \textbf{оди́н часы́} 是错的——要么说 \textbf{э́ти часы́}，要么换词。这类词的正确用法是「整词记 ＋ 配复数形容词」（\S4.3）。}

\subhdr{2.6　练习：判性、变复数、补代词}

\begin{enumerate}[label=\textbf{第 \arabic* 题}]
\item （У4 p11--p12）判性并写出代词：музе́й ｜ пе́сня ｜ мо́ре ｜ тетра́дь ｜ па́па
\item （У4 p13）下面四组里，哪两个是阴性？день ｜ ночь ｜ портфе́ль ｜ посте́ль
\item （У5 p6--p8）变复数：студе́нт ｜ уро́к ｜ слова́рь ｜ кни́га ｜ окно́ ｜ упражне́ние
\item （У7 p18）写出复数：челове́к ｜ ребёнок ｜ друг ｜ дом
\item （У5 p9、У7 p19）下面哪个词没有复数？молоко́ ｜ кни́га ｜ студе́нт；哪个只有复数？часы́ ｜ окно́ ｜ дом
\item （У4 p17--p18）用 он／она́／оно́ 替换：Где соба́ка? — \blank до́ма. ｜ Когда́ уро́к? — \blank за́втра. ｜ Мо́ре там? — \blank там.
\item （У5 p11--p14）按模型写一段描述：\textbf{Э́то го́род. Здесь проспе́кты, у́лицы, пло́щади …}（把括号里的名词全部变成复数）。
\end{enumerate}

\textbf{第 1 题}　关键一步：先看词尾落在 2.2 表哪一行；-ь 走四条归纳规律。答案：музе́й（阳，он）｜ пе́сня（阴，она́）｜ мо́ре（中，оно́）｜ тетра́дь（阴，она́，-ь 且词干末尾是 т，四条规律不符 → 查 2.3 右栏）｜ па́па（阳，он，特例）。

\textbf{第 2 题}　关键一步：只看 -ь 的两个词，查 2.3。答案：ночь（阴）、посте́ль（阴）；день 与 портфе́ль 是阳性。

\textbf{第 3 题}　关键一步：先判词干软硬。答案：студе́нты ｜ уро́ки（к 后写 и）｜ словари́ ｜ кни́ги（г 后写 и）｜ о́кна ｜ упражне́ния。

\textbf{第 4 题}　关键一步：这四个都在 2.5 的特殊表里。答案：лю́ди ｜ де́ти ｜ друзья́ ｜ дома́。

\textbf{第 5 题}　关键一步：只有单数／只有复数各查一张清单。答案：молоко́ 没有复数；часы́ 只有复数。

\textbf{第 6 题}　关键一步：соба́ка（-а → она́）、уро́к（辅音 → он）、мо́ре（-е → оно́）。答案：Она́ ｜ Он ｜ Оно́。

\textbf{第 7 题}　这道是自由写作题，课件没给答案，只给关键一步：\textbf{先把列表里的名词统一变复数}（\S2.4 的表），\textbf{再给每个复数配 -ые／-ие 形容词}（\S4.2），句子骨架照抄模型 \textbf{Э́то го́род. Здесь проспе́кты, у́лицы, пло́щади …}——不需要造新句子，只需要把名词换成复数形式。
